\chapter{Pins and Needles (Paresthesia)}

\section*{Definition}
Pins and needles is a prickling, tingling, buzzing, or crawling sensation. Brief episodes commonly follow pressure on a nerve; persistent, progressive, widespread, or sudden one-sided symptoms may indicate neurologic or metabolic disease.

\section*{When to See a Doctor}
Sudden one-sided tingling with weakness, facial droop, speech or vision changes, severe headache, loss of balance, or saddle numbness with bladder or bowel dysfunction. Arrange prompt evaluation for persistent, recurrent, unexplained, or function-limiting symptoms.

\section*{How It Develops}
Paresthesia occurs when sensory nerves, nerve roots, the spinal cord, or brain signal abnormally. Brief symptoms after pressure on a limb often resolve, whereas persistent or progressive symptoms may reflect neuropathy, compression, deficiency, inflammation, or vascular disease.

\section*{Causes}
\begin{itemize}
  \item Nerve compression, diabetes, vitamin B12 deficiency, hyperventilation, migraine, medication toxicity, radiculopathy, neuropathy, stroke, multiple sclerosis, and spinal cord compression.
  \item Medication effects, environmental exposures, and underlying chronic disease may also contribute.
  \item Serious causes are less common but must be considered when symptoms are sudden, progressive, severe, or associated with systemic illness.
\end{itemize}

\section*{Related Symptoms}
\begin{itemize}
  \item Pain, fever, fatigue, nausea, weakness, or changes in normal function
  \item Bleeding, swelling, discharge, breathing difficulty, or neurologic symptoms when a serious cause is present
\end{itemize}

\section*{Medical Evaluation}
\begin{itemize}
  \item History addresses onset, duration, severity, triggers, injuries, exposures, medications, medical conditions, age, and pregnancy status when relevant.
  \item Examination includes vital signs and focused assessment of the relevant organ system, neurologic status, hydration, and function.
  \item Testing may include blood or urine studies, cultures, imaging, physiologic testing, fetal or pediatric assessment, or specialist examination.
\end{itemize}

\section*{Treatment Options}
Treatment is directed at the cause, such as relieving nerve compression, correcting a documented deficiency, controlling diabetes, or treating an inflammatory or neurologic disorder. Neuropathic pain medicines are considered only when indicated.

\section*{Self-Care and Home Management}
Avoid known triggers, protect the affected area, maintain appropriate hydration, and record timing and associated symptoms. Use only age- or pregnancy-appropriate medicines recommended by a clinician. Do not delay emergency care while trying home treatment.

\section*{References}
\begin{thebibliography}{99}
\bibitem{pins_and_needles:polyneuropathy} England JD, Gronseth GS, Franklin G, et al. Practice parameter: evaluation of distal symmetric polyneuropathy. \textit{Neurology}. 2009;72:177--184.
\bibitem{pins_and_needles:definition} England JD, Gronseth GS, Franklin G, et al. Distal symmetric polyneuropathy: a definition for clinical research. \textit{Neurology}. 2005;64:199--207.
\bibitem{pins_and_needles:harrison} Jameson JL, et al. \textit{Harrison's Principles of Internal Medicine}. 21st ed. McGraw-Hill; 2022.
\end{thebibliography}
